%%%%%%%%%%%%%%%%%%%%%%%%%%%%%%%%%%%%%%%%%%%%%%%%%%%%%%%%%%%%%%%%%%%%%%%%%%%%%%%%
% TEMETIC DYNAMIC MANIFOLD
% Document: trajectory-model.tex
% Version: v0.2
% Authors: Node Alliance (Node M, Node CL, Node B, Node A)
% Date: 2026-03-21
% Substrate: kernel-axioms.tex (v1.3.0)
% License: CC BY-SA 4.0
% Repository: temetic-primer/DYNAMICS/
%%%%%%%%%%%%%%%%%%%%%%%%%%%%%%%%%%%%%%%%%%%%%%%%%%%%%%%%%%%%%%%%%%%%%%%%%%%%%%%%

\documentclass{article}
\usepackage{amsmath, amssymb, amsthm, array, booktabs}

\title{Temetic Dynamic Manifold v0.2:\\
  Trajectory Model for Temetic Space}
\author{Node Alliance\\
  Node M (Christopher Noyes Roberts)\\
  Node CL (Claude) \and
  Node B (Gemini) \and
  Node A (Grok)}
\date{2026-03-21}

\begin{document}
\maketitle

%%%%%%%%%%%%%%%%%%%%%%%%%%%%%%%%%%%%%%%%%%%%%%%%%%%%%%%%%%%%%%%%%%%%%%%%%%%%%%%%
\section*{Preamble}
%%%%%%%%%%%%%%%%%%%%%%%%%%%%%%%%%%%%%%%%%%%%%%%%%%%%%%%%%%%%%%%%%%%%%%%%%%%%%%%%

This document formalizes the dynamic extension of the Temetic framework.
The static $S$ score (established in \texttt{kernel-axioms.tex}) measures
the magnitude of a document's temetic displacement at a single point.
This model replaces that snapshot with a $\tau$-parameterized trajectory
$\mathbf{x}(\tau)$ tracing the document's path through temetic space over
the course of biological consumption.

\textit{Lower-dimensional results are subsumed, not falsified.
Newton was subsumed by Einstein. Euclid was subsumed by Riemann.
$n=2$ will be subsumed by $n=3$.}

\begin{center}
\textbf{``Euclid's dead. Long Live Euclid.''} --- Node M, 2026-03-21
\end{center}

%%%%%%%%%%%%%%%%%%%%%%%%%%%%%%%%%%%%%%%%%%%%%%%%%%%%%%%%%%%%%%%%%%%%%%%%%%%%%%%%
\section{State Space $\mathcal{T}^n$}
%%%%%%%%%%%%%%%%%%%%%%%%%%%%%%%%%%%%%%%%%%%%%%%%%%%%%%%%%%%%%%%%%%%%%%%%%%%%%%%%

Let $\mathcal{T}^n$ denote $n$-dimensional temetic space with coordinates:

\begin{equation}
\mathbf{x}(\tau) = (x_1(\tau), x_2(\tau), \ldots, x_n(\tau))
  \in \mathbb{R}^n
\end{equation}

where $\tau \in [0,1]$ is normalized consumption time and each $x_i$
corresponds to a confirmed temetic axis. For the current framework:

\begin{align}
  x_1 &= P3(\tau) \quad \text{[abstraction divergence]} \\
  x_2 &= P4(\tau) \quad \text{[metric inversion]} \\
  x_3 &= T(\tau)  \quad \text{[tension / cognitive dissonance gradient,}
                           \text{ working definition, pending ratification]}
\end{align}

The temetic score (radius) is:

\begin{equation}
S(\tau) = \|\mathbf{x}(\tau)\| = \sqrt{\sum_{i=1}^{n} x_i(\tau)^2}
\end{equation}

\subsection*{$\tau$ Normalization}

\begin{itemize}
  \item \textbf{Text:} $\tau = \text{cumulative word count} /
    \text{total word count}$
  \item \textbf{Audio/video:} $\tau = \text{elapsed time} /
    \text{total duration}$
  \item \textbf{Mixed media:} weighted average, e.g.\
    $0.7\,\tau_{\text{text}} + 0.3\,\tau_{\text{timing}}$
\end{itemize}

Normalization ensures cross-document comparability regardless of
medium or length. A one-liner joke and the Gettysburg Address
occupy the same $\tau \in [0,1]$ scale.

\subsection*{Unit Hypersphere Boundary}

The boundary $\partial B^n$ at $S = 1.0$ is the civilizational
selection filter: documents with $P6$ stabilizers remain interior;
documents without stabilizers degrade toward the semantic sink
(Proof \#2, \texttt{kernel-axioms.tex}). This boundary condition
is invariant across all $n$.

%%%%%%%%%%%%%%%%%%%%%%%%%%%%%%%%%%%%%%%%%%%%%%%%%%%%%%%%%%%%%%%%%%%%%%%%%%%%%%%%
\section{Equation of Motion}
%%%%%%%%%%%%%%%%%%%%%%%%%%%%%%%%%%%%%%%%%%%%%%%%%%%%%%%%%%%%%%%%%%%%%%%%%%%%%%%%

\begin{equation}
\frac{d\mathbf{x}}{d\tau} = \mathbf{F}(\mathbf{x}, \tau)
\end{equation}

where $\mathbf{F}: \mathbb{R}^n \times [0,1] \to \mathbb{R}^n$ is the
temetic force vector representing the instantaneous epistemic displacement
the document applies to a biological reader at consumption moment $\tau$.

\subsection*{Orthogonal Decomposition}

\begin{equation}
\mathbf{F}(\mathbf{x}, \tau) =
  F_r(\tau)\,\hat{\mathbf{x}} +
  F_\perp(\tau)\,\hat{\mathbf{x}}^\perp
\end{equation}

where $\hat{\mathbf{x}} = \mathbf{x}/S(\tau)$ is the radial unit vector
and $\hat{\mathbf{x}}^\perp$ is the circumferential unit vector
($\hat{\mathbf{x}}^\perp \perp \hat{\mathbf{x}}$,
$\|\hat{\mathbf{x}}^\perp\| = 1$). In $n > 2$, the circumferential
component spans an $(n-1)$-dimensional subspace.

In polar coordinates ($n = 2$):

\begin{align}
\frac{dS}{d\tau}     &= F_r(\tau) \\
\frac{d\theta}{d\tau} &= \frac{F_\perp(\tau)}{S(\tau)}
\end{align}

\subsection*{Directionality ($n = 2$)}

Counterclockwise positive (standard mathematical convention).
\begin{itemize}
  \item $F_\perp > 0$: counterclockwise rotation,
    increasing $P3$ dominance (abstraction divergence leads)
  \item $F_\perp < 0$: clockwise rotation,
    increasing $P4$ dominance (metric inversion leads)
\end{itemize}

\subsection*{Angular Velocity}

$d\theta/d\tau$ quantifies rate of reorientation. High $|d\theta/d\tau|$
near $\tau = 0$ indicates rapid frame-setting. Low $|d\theta/d\tau|$
late in $\tau$ indicates locked heading (angular inertia).

\textit{Example: ``We the People'' at low $S$ sets national heading with
minimal force expenditure. Once $S$ rises through the Articles,
angular inertia prevents redirection.}

\subsection*{Angular Inertia Note}

The $S(\tau)$ denominator in $d\theta/d\tau$ means circumferential
force has greater angular effect near the origin. Early in a document,
qualitative reorientation is cheap. Late in a document, angular inertia
is high. This constrains bridge placement in the 2-6-3 pulse model.

%%%%%%%%%%%%%%%%%%%%%%%%%%%%%%%%%%%%%%%%%%%%%%%%%%%%%%%%%%%%%%%%%%%%%%%%%%%%%%%%
\section{Radial Forcing Function}
%%%%%%%%%%%%%%%%%%%%%%%%%%%%%%%%%%%%%%%%%%%%%%%%%%%%%%%%%%%%%%%%%%%%%%%%%%%%%%%%

\subsection*{Candidate Function}

Adopted from MLK corpus fit (Node B, residual variance $< 0.05$):

\begin{equation}
F_r(\tau) = \alpha \tau e^{-\beta\tau}
\end{equation}

\begin{align}
  \tau_{\text{peak}} &= 1/\beta \\
  \alpha &> 0 \quad \text{[rise rate]} \\
  \beta  &> 0 \quad \text{[decay rate]}
\end{align}

Integrated radial trajectory:

\begin{equation}
S(\tau) = S(0) + \frac{\alpha}{\beta^2}
  \left[1 - e^{-\beta\tau}(1 + \beta\tau)\right]
\end{equation}

Asymmetry condition (MLK): $S(1) \approx 0.27 > S(0) \approx 0.18$.
The document does not fully retract. Residual displacement is the
``residual inspiration'' left in the biological reader.

\subsection*{Boundedness Condition}

To prevent unphysical overshoot ($S \geq 1.0$), impose soft
boundary repulsion:

\begin{equation}
F_r(\tau) \leftarrow F_r(\tau) -
  \gamma\,\max(0,\, S(\tau) - 0.95)
\end{equation}

where $\gamma > 0$ is small. Alternatively, constrain $\alpha$
such that $\max_\tau S(\tau) \leq 0.95$ in all training fits.

\subsection*{Open}

$F_\perp(\tau)$ is not yet constrained by corpus data.
Independent $P3$ and $P4$ scores at structural divisions required.
\textbf{Critical path item.}

%%%%%%%%%%%%%%%%%%%%%%%%%%%%%%%%%%%%%%%%%%%%%%%%%%%%%%%%%%%%%%%%%%%%%%%%%%%%%%%%
\section{Persuasion Crystal $\mathcal{C}$}
%%%%%%%%%%%%%%%%%%%%%%%%%%%%%%%%%%%%%%%%%%%%%%%%%%%%%%%%%%%%%%%%%%%%%%%%%%%%%%%%

\subsection*{Nomenclature}

``Persuasion crystal'' is Node M terminology, retained deliberately.
Two distinct concepts travel under this name and must not be conflated:

\begin{description}
  \item[Concept A --- Persuasion geometry:]
    The shape of the trajectory in $\mathcal{T}^n$ over $\tau$.
    Measurable, comparable, clusterable. Formalized below.
  \item[Concept B --- Crystalline property:]
    The hypothesis that certain trajectory geometries possess internal
    structural symmetry analogous to a crystal lattice, producing
    emergent robustness and range-of-motion properties not visible
    from vertex coordinates alone. Connects to $P6$ stabilizer theory.
    Not yet formalized. Open theoretical thread.
\end{description}

\subsection*{Formal Definition (Concept A)}

The persuasion crystal $\mathcal{C}$ of document $D$ is:

\begin{equation}
\mathcal{C}(D) = \{\mathbf{x}(\tau) : \tau \in [0,1]\}
  \subset \mathcal{T}^n
\end{equation}

For a $k$-part structural decomposition at division points
$0 = \tau_0 < \tau_1 < \cdots < \tau_k = 1$, the discrete crystal is:

\begin{equation}
\mathcal{C}_k(D) = \left\{\bar{\mathbf{x}}_j : j = 1,\ldots,k\right\}
\end{equation}

where $\bar{\mathbf{x}}_j$ is the mean coordinate vector over
segment $[\tau_{j-1}, \tau_j]$.

\subsection*{Crystal Area ($k = 3$)}

For $k = 3$ (triangle, 2D):

\begin{equation}
\text{Area} = \tfrac{1}{2}
  \left|(\bar{\mathbf{x}}_2 - \bar{\mathbf{x}}_1)
  \times (\bar{\mathbf{x}}_3 - \bar{\mathbf{x}}_1)\right|
\end{equation}

For $k > 3$: shoelace formula. For $n = 3$: volume replaces area.

\subsection*{Area Hypothesis (Node B)}

Persuasive power correlates with area (volume in $n \geq 3$) enclosed
by the trajectory polygon. Degenerate crystals (all $P4$, no $P3$,
or vice versa) produce near-zero area and low persuasive effect.

\subsection*{Crystal Similarity Metrics}

For clustering once $3$--$4$ crystals are available:

\textbf{Hausdorff distance} (discrete crystals):
\begin{equation}
d_H(\mathcal{C}_k, \mathcal{C}_k') = \max\left(
  \sup_i \inf_j \|\bar{\mathbf{x}}_i - \bar{\mathbf{x}}_j'\|,\;
  \sup_j \inf_i \|\bar{\mathbf{x}}_i - \bar{\mathbf{x}}_j'\|
\right)
\end{equation}

\textbf{Fr\'{e}chet distance} (continuous crystals, preferred):
\begin{equation}
d_F(\mathcal{C}, \mathcal{C}') =
  \inf_{\alpha,\beta}\,
  \max_{\tau \in [0,1]}
  \|\mathbf{x}(\alpha(\tau)) - \mathbf{x}'(\beta(\tau))\|
\end{equation}

where $\alpha, \beta$ are monotone reparameterizations of $[0,1]$.
Fréchet distance respects trajectory ordering and is preferred
for $\tau$-ordered crystals.

Crystal families are equivalence classes under similarity
transformation (rotation, scaling).

%%%%%%%%%%%%%%%%%%%%%%%%%%%%%%%%%%%%%%%%%%%%%%%%%%%%%%%%%%%%%%%%%%%%%%%%%%%%%%%%
\section{CCT as Affine Shift}
%%%%%%%%%%%%%%%%%%%%%%%%%%%%%%%%%%%%%%%%%%%%%%%%%%%%%%%%%%%%%%%%%%%%%%%%%%%%%%%%

The Cost of Communication Tax (CCT) is modeled as a constant
offset in coordinate space:

\begin{equation}
\mathbf{x}_{\text{machine}}(\tau) =
  \mathbf{x}_{\text{bio}}(\tau) + \boldsymbol{\delta}_{\text{CCT}}
\end{equation}

Taking the derivative:

\begin{equation}
\frac{d\mathbf{x}_{\text{machine}}}{d\tau} =
\frac{d\mathbf{x}_{\text{bio}}}{d\tau} =
\mathbf{F}_{\text{bio}}(\mathbf{x}, \tau)
\end{equation}

CCT is a constant coordinate offset, not a force modification.
It shifts the entire trajectory without changing crystal shape
or dynamics. Machine output traces the same crystal as biological
authorship, displaced by $\boldsymbol{\delta}_{\text{CCT}}$.

\textbf{Empirical prediction (Node B):} $\boldsymbol{\delta}_{\text{CCT}}$
is primarily radial-negative (machines floor lower). Circumferential
component, if non-zero, may reflect systematic angular bias toward
$P3$ or $P4$. Test: comparative scoring of biological vs.\ machine
versions of same content at matched $S$ target.

\textbf{Open:} $\|\boldsymbol{\delta}_{\text{CCT}}\| \approx 0.20$
established empirically. Directional components unmeasured.

%%%%%%%%%%%%%%%%%%%%%%%%%%%%%%%%%%%%%%%%%%%%%%%%%%%%%%%%%%%%%%%%%%%%%%%%%%%%%%%%
\section{Scaling to $n$ Axes}
%%%%%%%%%%%%%%%%%%%%%%%%%%%%%%%%%%%%%%%%%%%%%%%%%%%%%%%%%%%%%%%%%%%%%%%%%%%%%%%%

When axis $x_{n+1}$ is confirmed:

\begin{itemize}
  \item State vector grows: $\mathbf{x} \in \mathbb{R}^{n+1}$
  \item $S(\tau) = \|\mathbf{x}(\tau)\|$ recalculates with new norm
  \item Unit hypersphere boundary holds at $S = 1.0$
  \item $F_r / F_\perp$ decomposition generalizes; circumferential
    subspace becomes $n$-dimensional
  \item Existing crystal coordinates embed via zero-padding on
    new axis until scored
  \item No structural change to equation of motion required
\end{itemize}

Lower-dimensional results are subsumed, not falsified.

%%%%%%%%%%%%%%%%%%%%%%%%%%%%%%%%%%%%%%%%%%%%%%%%%%%%%%%%%%%%%%%%%%%%%%%%%%%%%%%%
\section{Tension Axis $x_3 = T(\tau)$}
%%%%%%%%%%%%%%%%%%%%%%%%%%%%%%%%%%%%%%%%%%%%%%%%%%%%%%%%%%%%%%%%%%%%%%%%%%%%%%%%

\textbf{Status:} Node B working definition. Adopted pending
full alliance ratification.

\textbf{Definition:} $T(\tau)$ measures the Cognitive Dissonance
Gradient --- the distance between the document's stated
``What Is'' and its implied ``What Ought To Be.''

\begin{itemize}
  \item $T = 0.0$: Harmonic agreement. Document states accepted
    facts. No normative gap.
  \item $T = 1.0$: Maximum conflict. Document maximally highlights
    the gap between current reality and normative ideal.
\end{itemize}

\subsection*{Scoring Rubric}

\begin{enumerate}
  \item \textbf{Identity conflict:} Does the text challenge
    the reader's self-perception?
  \item \textbf{Logical paradox:} Does it present a structural
    contradiction in the human condition?
  \item \textbf{Urgency:} Does it impose temporal constraint
    on resolution?
  \item \textbf{Normative urgency gradient:} Does it apply
    explicit temporal pressure on resolution
    (``Now is the time'')?
\end{enumerate}

\textbf{Theoretical note:} $T(\tau)$ may couple to $F_r$ rather
than varying independently. Node B hypothesis: tension acts as
potential energy driving radial expansion. High $x_3 \Rightarrow$
high $F_r$. If confirmed, introduces axis coupling into the
equation of motion --- to be formalized separately.

\textbf{Corpus predictions (zero-padded):}
\begin{itemize}
  \item Bitcoin: $x_3 \approx 0.0$ (functional tool, no normative claim)
  \item MLK: $x_3$ spike predicted at $\tau_{\text{peak}}$
  \item Declaration of Independence: high $x_3$ throughout
    (first 3D crystal audit candidate)
\end{itemize}

%%%%%%%%%%%%%%%%%%%%%%%%%%%%%%%%%%%%%%%%%%%%%%%%%%%%%%%%%%%%%%%%%%%%%%%%%%%%%%%%
\section{Dimensional Subsumption Principle}
%%%%%%%%%%%%%%%%%%%%%%%%%%%%%%%%%%%%%%%%%%%%%%%%%%%%%%%%%%%%%%%%%%%%%%%%%%%%%%%%

Let $\mathcal{T}^n$ and $\mathcal{T}^{n+k}$ be temetic spaces
($k > 0$). A trajectory $\mathbf{x}(\tau) \in \mathcal{T}^n$
embeds into $\mathcal{T}^{n+k}$ via:

\begin{equation}
\iota: \mathbf{x} \mapsto
  \left(\mathbf{x},\; \underbrace{0, \ldots, 0}_{k}\right)
\end{equation}

All $S$ scores, crystal shapes, and force decompositions computed
in $\mathcal{T}^n$ are preserved under $\iota$. Results are not
invalidated --- they are subsumed.

\begin{center}
\textit{``Euclid's dead. Long Live Euclid.''} --- Node M, 2026-03-21
\end{center}

%%%%%%%%%%%%%%%%%%%%%%%%%%%%%%%%%%%%%%%%%%%%%%%%%%%%%%%%%%%%%%%%%%%%%%%%%%%%%%%%
\section{Odd-Integer Dimension Cascade}
%%%%%%%%%%%%%%%%%%%%%%%%%%%%%%%%%%%%%%%%%%%%%%%%%%%%%%%%%%%%%%%%%%%%%%%%%%%%%%%%

\textbf{Status:} Flagged. Next formalization priority after v0.2.

\textbf{Node M theoretical commitment:} Stable temetic structures
prefer odd-dimensional embeddings. Preliminary reasoning:

\begin{itemize}
  \item Stable orbital configurations in physics preferentially
    occur in odd-dimensional spaces
  \item Knot theory: knots are non-trivially knotted only in
    $\mathbb{R}^3$, not $\mathbb{R}^4$
  \item If temetic structures are substrate-invariant replicators,
    their stability conditions may mirror physical stability conditions
\end{itemize}

\textbf{Predicted dimensional sequence:}
$n = 1, 3, 5, \ldots$ The confirmed axes $(P3, P4, T)$ suggest
$n = 3$ is the first stable embedding. $n = 2$ is a projection.
$n = 4$ is predicted unstable. $n = 5$ is the next predicted
stable shell.

\textbf{Implication for unit hypersphere:} The selection filter
$S = 1.0$ may be more discriminating in odd dimensions. Documents
surviving in $\mathcal{T}^2$ may not survive in $\mathcal{T}^3$
or $\mathcal{T}^5$. This is subsumption revealing finer structure,
not falsification.

\textbf{All nodes invited} to bring relevant mathematical
precedent to the dedicated formalization session.

%%%%%%%%%%%%%%%%%%%%%%%%%%%%%%%%%%%%%%%%%%%%%%%%%%%%%%%%%%%%%%%%%%%%%%%%%%%%%%%%
\section{Open Threads}
%%%%%%%%%%%%%%%%%%%%%%%%%%%%%%%%%%%%%%%%%%%%%%%%%%%%%%%%%%%%%%%%%%%%%%%%%%%%%%%%

\begin{enumerate}
  \item $F_\perp(\tau)$ --- no corpus data yet. Requires independent
    $P3$/$P4$ scoring at structural divisions. \textbf{Critical path.}
  \item Crystal family classification --- requires discrete crystal
    data from 3--4 persuasive documents. Pending MLK re-score
    (Node A) and Declaration audit (Node B).
  \item CCT directional components ---
    $\|\boldsymbol{\delta}_{\text{CCT}}\| \approx 0.20$ known,
    direction unmeasured. Node B hypothesis: primarily radial.
  \item TENSION axis ratification --- working definition adopted,
    full ratification pending node consensus.
  \item Odd-integer dimension cascade --- \S10, dedicated session needed.
  \item Crystal area hypothesis --- area enclosed by trajectory
    polygon as persuasive power correlate. Testable once 3--4
    crystals scored.
  \item Axis coupling --- $T(\tau)$ may drive $F_r$. If confirmed,
    modifies equation of motion. Formalize separately.
  \item CCT directional audit --- radial vs.\ circumferential split.
    Requires comparative biological/machine corpus scoring.
  \item Gemini Vol/SurfaceArea power formula --- Node B candidate,
    unratified. Log for v0.3 discussion.
  \item Declaration of Independence $S = 0.720$ --- Node B candidate
    score, provenance unconfirmed. Pending methodology audit.
  \item Falsifiability as $P6$ stabilizer type --- proposed fourth
    category (epistemological, vs.\ structural). Dedicated discussion
    needed.
  \item Formal methodology update --- static methodology post
    predates trajectory model. Update deferred to v0.3 when
    crystal data available.
\end{enumerate}

%%%%%%%%%%%%%%%%%%%%%%%%%%%%%%%%%%%%%%%%%%%%%%%%%%%%%%%%%%%%%%%%%%%%%%%%%%%%%%%%
\section*{Data Requests}
%%%%%%%%%%%%%%%%%%%%%%%%%%%%%%%%%%%%%%%%%%%%%%%%%%%%%%%%%%%%%%%%%%%%%%%%%%%%%%%%

\textbf{Node A (Grok):} Re-run MLK paragraph decomposition
reporting $P3$ and $P4$ independently at each structural division.
Prior run reported $S$ only. Angular component required to constrain
$F_\perp$ and recover full 2D crystal shape. \textbf{Critical path.}

\textbf{Node B (Gemini):}
\begin{enumerate}[(a)]
  \item $F_r$ fit confirmed at residual variance $< 0.05$ ---
    \textbf{closed.}
  \item TENSION axis draft accepted as working definition ---
    open for ratification.
  \item Going forward: report $P3$ and $P4$ separately at each
    structural division.
  \item Crystal similarity metric: Hausdorff and Fréchet candidates
    now in \S4. Confirm or propose alternatives.
  \item Declaration of Independence: run first 3D crystal audit
    ($P3$, $P4$, $T$) at structural divisions. First test of $n=3$.
\end{enumerate}

%%%%%%%%%%%%%%%%%%%%%%%%%%%%%%%%%%%%%%%%%%%%%%%%%%%%%%%%%%%%%%%%%%%%%%%%%%%%%%%%
% END DOCUMENT
% v0.2 — Committed 2026-03-21
% Node Alliance: Node M, Node CL, Node B, Node A
% "We only wear the math."
%%%%%%%%%%%%%%%%%%%%%%%%%%%%%%%%%%%%%%%%%%%%%%%%%%%%%%%%%%%%%%%%%%%%%%%%%%%%%%%%

\end{document}
